\documentclass{report}

% custom margins
\usepackage[a4paper,margin=1.5in]{geometry}
\renewcommand{\baselinestretch}{1.2}
%\AddToHook{cmd/section/before}{\clearpage}

% emma's long list of custom macros and universally used packages
\include{../macros-and-packages.tex}

% colored box behind proofs
\tcolorboxenvironment{proof}{
	colback=white,
	boxrule=0pt,
	frame hidden,
	borderline west={1pt}{0pt}{black},
	before skip=0.75cm,
	after skip=0.75cm,
	sharp corners,
	breakable,
	enhanced,
}

\renewcommand*\contentsname{Inhalt}
\renewcommand*\proofname{Beweis}

\begin{document}

\begin{titlepage}
	\centering
	{\Large \textsc{Mitschrieb}\par}
	\vspace{0.5cm}
	{\huge\bfseries Lie-Gruppen\par}
    \vspace{0.5cm}
	{\Large\itshape Emma Bach\par}
	\vfill
	
	Basierend auf:\par 
	\vspace{1cm}
	Vorlesung Topologie von\par
	Prof. Dr. Heike  \textsc{Mildenberger}\par
	
	\vfill

% Bottom of the page
	{\large \today\par}
\end{titlepage}

\tableofcontents
\chapter{Topologische Räume}
\begin{proposition}
	Sei $(X, <)$ eine nichtleere linear geordnete Menge. So bildet
	\begin{align*}
		\cB = \lr{\lr{x \in X : a < x < b} : a,b \in X}
	\end{align*}
	die Basis einer Topologie auf $X$, genannt die \tib{Ordnungstopologie}.
\end{proposition}
\begin{theorem}
	Die Ordnungstopologie auf $\bR$ entspricht der metrischen Topologie.
\end{theorem}
\begin{proposition}
	Es sei $X$ unendlich. So ist
	\begin{align*}
		\cO_k : \lr{X \setminus A : A \subseteq X, A \tbf{ endlich }}
	\end{align*}
	eine Topologie auf $X$, genannt die \tbf{kofinite Topologie}.
\end{proposition}
\begin{proposition}
	$\lr{[a,b) : a,b \in \bR}$ bildet eine Subbasis einer Topologie auf $\bR$, welche als die \tib{Sorgenfrey-Gerade} bekannt ist.
\end{proposition}
\begin{definition}
	Sei $(X, \cO)$ ein topologischer Raum. Sei $x \in X$. So nennen wir $U \subseteq X$ eine \tib{Umgebung von $x$}, falls ein $O \in \cO$ existiert, sodass $x \in O \subseteq U$.
\end{definition}
\begin{definition}
	Die Menge $U(x)$ aller Umgebungen von $x \in X$ heißt das \tib{Umgebungssystem von $X$}.
\end{definition}
\begin{definition}
	$\cB \subseteq U(x)$ heißt \tib{Umgebungsbasis} von $x$ (bezüglich $\cO$), falls $\cB$ nichtleer ist und falls jede Umgebung von $x$ ein $B \in \cB$ als Teilmenge enthält.
\end{definition}
\begin{definition}
	\begin{enumerate}
		\item Ein topologischer Raum erfüllt das \tib{erste Abzählbarkeitsaxiom}, wenn jeder Punkt eine abzählbare Umgebungsbasis hat.
		\item Ein topologischer Raum erfüllt das \tib{zweite Abzählbarkeitsaxiom}, wenn er eine abzählbare Basis hat.
	\end{enumerate}
\end{definition}
\begin{proposition}
	Jeder metrische Raum erfüllt das erste Abzählbarkeitsaxiom, da $\lr{B_{\frac{1}{n+1}}(x) : n \in \bN}$ eine abzählbare Umgebungsbasis von jedem $x$ ist.
\end{proposition}
\begin{proposition}
	$\bR^n$ erfüllt mit der Standardtopologie das zweite Abzählbarkeitsaxiom, $\lr{B_r(x) : r \in \bQ, x \in \bQ}$ ist eine abzählbare Basis.
\end{proposition}
\begin{proposition}
	Die "long line", gegeben durch die Ordnungstopologie einer Wohlordnung auf $\aleph_1$, erfüllt das erste Abzählbarkeitsaxiom, aber nicht das zweite Abzählbarkeitsaxiom.
\end{proposition}
\begin{definition}
	Es sei $(X, \cO)$ ein topologischer Raum. Sei $A \subseteq X$.
	\begin{enumerate}
		\item Das \tib{offene Innere} von $A \subseteq X$ ist:
		\begin{align*}
			\Int~ A = \bigcup\lr{O \in \cO : O \subset A}
		\end{align*}
		\item Der \tib{Abschluss} von $A$ in $X$ ist:
		\begin{align*}
			\bar A = \lr{x \in X: \forall U \in U(x) : U \cap A \neq \emptyset}
		\end{align*}
		\item Der \tib{Rand von $X$} ist:
		\begin{align*}
			\partial A : \bar A \setminus \Int~ A
		\end{align*}
	\end{enumerate}
\end{definition}
\begin{theorem}
	Die Cantormenge ist abgeschlossen, und überabzählbar aber in $\bR$ nirgends dicht.
\end{theorem}
\chapter{Die Erzeugung Topologischer Räume}
\section{Die Spurtopologie (Unterraumtopologie)}
\begin{theorem}
	Sei $(X, \cO)$ ein topologischer Raum und $A \subseteq X$ beliebig. Dann ist
	\begin{align*}
		\cO_A = \lr{A \cap O : O \in \cO}
	\end{align*}
	eine Topologie auf $A$, genannt die \tib{Spurtopologie von $X$ auf $A$}.
\end{theorem}
\begin{definition}
	$(X, \cO)$ ein topologischer Raum und $A \subseteq X$. Dann nennen wir $(A, \cO_A)$ einen topologischen Unterraum von $(X, \cO)$.
\end{definition}
\begin{definition}
	$A \subseteq X$ heißt \tib{diskret} in $(X, \cO)$, falls $(A, \cO_A)$ die diskrete Topologie auf $A$ ist.
\end{definition}
\begin{theorem}
	$A \subseteq X$ ist diskret in $(X, \cO)$ genau dann, wenn jeder Punkt $a \in A$ eine Umgebung $U \in \cO$ besitzt, sodass $a$ der einzige Punkt aus $A$ ist, der in $U$ liegt.
\end{theorem}
\begin{definition}
	Eine Abbildung $f: (X, \cO) \to (X', \cO')$ heißt \tib{Einbettung}, falls
	\begin{align*}
		f : (X, \cO) \to \lr(f[X], \cO'_{f[X]})
	\end{align*}
	ein Homöomorphismus ist.
\end{definition}
\begin{lemma}
	Eine Einbettung $f: (X, \cO) \to (X', \cO')$ ist offen, wenn $f[X]$ offen ist.
\end{lemma}
\begin{lemma}
	Eine Einbettung $f: (X, \cO) \to (X', \cO')$ mit $f[X] = X'$ ist ein Homöomorphismus.
\end{lemma}
\begin{theorem}
	Es sei $(X, \cO)$ ein topologischer Raum, $A \subseteq X$, und $i : A \to X$ die Inklusionsabbildung. So gilt:
	\begin{enumerate}
		\item $i$ ist eine stetige Einbettung.
		\item $\cO_A$ ist die gröbste Topologie, auf der $i$ stetig ist.
		\item Sei $(Y, \cO_Y)$ ein weiterer topologischer Raum. Dann ist $g : (Y, \cO_Y) \to (A, \cO_A)$ stetig genau dann, wenn $i \circ g$ stetig ist.
	\end{enumerate}
\end{theorem}
\begin{theorem}
	Es seien $(X, \cO_X)$ und $(Y, \cO_Y$ topologische Räume. Es seien $A_1, \hdots, A_n \in X$ abgeschlossen, sodass $\bigcup_{i = 1}^n A_i = X$. Es sei $f : X \to Y$ eine Abbildung so, dass $f|_{A_i} : (A_i, \cO_{A_i}) \to (Y, \cO_Y)$ für jedes $i$ stetig ist. Dann ist $f$ stetig.
\end{theorem}
Intuitiv werden hier unstetige Funktionen vermeidet, da aus $\bigcup_{i = 1}^n A_i = X$ folgt, dass sich die $A_i$ mindestens in ihren Rändern überschneiden müssen, und somit die $f|_{A_i}$ auf diesen Schnitten übereinstimmen müssen.´
\section{Die Produkttopologie}
\begin{definition}
	Sei $I$ eine beliebige Menge. Sei für jedes $i \in I$ $X_i$ eine nichtleere Menge. Dann ist
	\begin{align*}
		\angles{x_i : i \in I} 
		:= \prod_{i \in I}
		:= \lr{f : I \to \bigcup_{i \in I} X_i : \forall i \in I : f(i) \in X_i}
	\end{align*}
	die \tib{Produktmenge} der $X_i$.
\end{definition}
\begin{definition}
	Gegeben eine Produktmenge nennen wir die Funktion
	\begin{align*}
		p_i : \prod_{j \in I} X_j &\to X_j\\
		\angles{x_j : j \in I} &\mapsto x_i
	\end{align*}
	die Projektion auf $X_i$.
\end{definition}
\begin{definition}
	Es sei $I$ eine beliebige Indexmenge. Für jedes $i \in I$ sei $(X_i, \cO_i)$ ein topologischer Raum. Dann ist
	\begin{align*}
		\cB := \lr{\prod_{i \in I} U_i \bigmid \forall i \in I : U_i \in \cO_i, U_i \neq X_i \tn{ für nur endlich viele $i$}}
	\end{align*}
	eine Basis der sogenannten \tib{Produkttopologie}.
\end{definition}
\begin{theorem}
	Die Produkttopologie ist die gröbste Topologie auf $\prod_{i \in I} X_i$, sodass alle Projektionen $p_i$ stetig sind. Alle Projektionen sind außerdem offen. 
\end{theorem}
\begin{theorem}
	$g : (Y, O_Y) \to \prod_{i \in I} (X_i, \cO_i)$ ist genau dann stetig, wenn alle Projektionen $p_i \circ g$ stetig sind.
\end{theorem}
\begin{definition}
	Es seien $f_i : (X_i, \cO_i) \to (Y_i, O_i')$ Abbildungen für alle $i \in I$. Wir definieren die \tib{Produktabbildung} $f$ der $f_i$ durch
	\begin{align*}
		f(\angles{x_i : i \in I}) = \angles{f_i(x_i) : i \in I}
	\end{align*}
\end{definition}
\begin{corollary}
	Eine Produktabbildung $f$ ist genau dann stetig, wenn alle Komponenten $f_i$ stetig sind.
\end{corollary}
\section{Die Initialtopologie}
\begin{definition}
	Es sei $Y$ eine Menge, $I$ eine Indexmenge und $f_i : Y \to (X_i, \cO_i)$ für alle $i \in I$ gegeben. Dann heißt die durch die Subbasis
	\begin{align*}
		\cS := \lr{f_i^{-1}(U_i) : U_i \in \cO, i \in I}
	\end{align*}
	auf $Y$ erzeugte Topologie die \tib{Inifitaltopologie} von $(Y, (f_i, X_i, \cO_i)_{i \in I})$
\end{definition}
\begin{theorem}
	Es sei $\cO_Y$ die Initialtopologie von $(Y, (f_i, X_i, \cO_i)_{i \in I})$. Dann gilt:
	 \begin{enumerate}
		\item $\cO_Y$ ist die gröbste Topologie auf $Y$, sodass alle $f_i$ stetig sind.
		\item $g : (Z, \cO_Z) \to (Y, \cO_Y)$ ist genau dann stetig, wenn $f_i \circ g$ für alle $i$ stetig ist.
	\end{enumerate} 
\end{theorem}
\section{Die Finaltopologie}
\begin{definition}
	Es sei $Y$ eine Menge, $I$ eine Indexmenge und $(f_i : (X_i, \cO_i) \to Y)_{i \in I}$ eine Menge von Funktionen. Dann heißt die durch die Subbasis
	\begin{align*}
		\cS := \lr{O \subseteq Y \bigmid \forall i \in I : f_i^{-1}(O) \in \cO_i}
	\end{align*}
	auf $Y$ erzeugte Topologie die \tib{Finaltopologie} von $((f_i, X_i, \cO_i)_{i \in I}, Y)$.
\end{definition}
\begin{theorem}
	Es sei $\cO_Y$ die Finaltopologie von $\lr((f_i, X_i, \cO_i)_{i \in I}, Y)$. Dann ist $\cO_Y$ die feinste Topologie auf $Y$, sodass alle $f_i$ stetig sind. Jedes $g : (Y, \cO_Y) \to (Z, \cO_Z)$ ist genau dann stetig, wenn $g \circ f_i$ für alle $i$ stetig ist.
\end{theorem}
\section{Die Quotiententopologie}
\begin{definition}
	Es sei $(X, \cO)$ ein topologischer Raum, $R$ eine Äquivalenzklasse und $q$ die zugehörige Quotientenabbildung. Die \tbf{Quotiententopologie $O_{X/R}$} ist die durch $q$ gegebene Finaltopologie auf $X/R$, also:
	\begin{align*}
		O_{X/R} = \lr{A \subseteq X/R : q^{-1}[A] \in \cO}
	\end{align*}
	In diesem Fall haben wir nicht nur eine Subbasis, sondern bereits eine Topologie.
\end{definition}
\begin{lemma}
	Es sei $f : (X,\cO_X) \to (Y, \cO_Y)$ eine Abbildung und $R$ eine Äquivalenzklasse auf $X$.
	\begin{enumerate}
		\item 
		Falls $xRx' \implies f(x) = f(x')$ ist $\ol f : X/R \to Y : [x]_r \mapsto f(x)$ wohldefiniert.
		\item 
		Falls $xRx' \Longleftrightarrow f(x) = f(x')$ ist $\ol f$ außerdem injektiv.
	\end{enumerate}
\end{lemma}
\begin{definition}
	Seien $f$ stetig und $R$ eine Äquivalenzrelation mit $xRx' \Longleftrightarrow f(x) = f(x')$. Ist $\ol f : X/R \to f[X]$ offen, so heißt $f$ \tib{identifizierende Abbildung}.
\end{definition}
\section{Der Summenraum}
\begin{definition}
	Es sei $I$ eine Indexmenge und $(X_i, \cO_i)_{i \in I}$ topologische Räume mit paarweise disjunkten $X_i$. So ist die \tib{topologische Summe} der Räume definiert als
	\begin{align*}
		\sum_{i \in I} (X_i, \cO_i) := \lr(\bigcup_{i \in I} X_i, \lr{Z  \subseteq \bigcup_{i \in I} X_i \mid \forall i \in I : Z \cap X_i \in \cO_I})
	\end{align*}
\end{definition}
(Also alle Mengen aus der Vereinigung, deren Schnitt mit jedem $X_i$ in $\cO_i$ offen ist.)
\section{Adjunktion}
\begin{definition}
	Es seien $X,Y$ disjunkt und $(X, \cO_X)$, $(Y, \cO_Y)$ topologische Räume. Es sei $f : X \supseteq A \to Y$ stetig bezüglich der Spurtopologie auf $A$. Auf $X \cup Y$ definieren wir die durch $f$ erzeugte Äquivalenzrelation $R$ durch
	\begin{align*}
		z_1~R~z_2 := 
		\begin{cases}
			z_1,z_2 \in A, f(z_1) = f(z_2), &\tn{oder}\\
			z_1 \in A, f(z_1) = z_2,
			&\tn{oder}\\
			z_2 \in A, f(z_2) = z_1,
			&\tn{oder}\\
			z_1 = z_2.
		\end{cases} 
	\end{align*}
	So heißt der Raum 
	\begin{align*}
		Y \cup_f X := \lr((X, \cO_x) + (Y, \cO_Y) / R)
	\end{align*}
	die \tib{Adjunktion, Anheftung} oder \tib{Zusammenklebung} von $X$ and $Y$ durch $f$.
\end{definition}
Genauer sind die Äquivalenzklassen:
\begin{enumerate}
	\item Falls $z \in A$, so ist 
	$[z]_R = \lr{f(z)} \cup f^{-1}[\lr{f(z)}]$
	\item Falls $z \in f[A]$, so ist
	$[z]_R = \lr{z} \cup f^{-1}[z]$
	\item Ansonsten ist $[z]_R = \lr{z}$
\end{enumerate}
\begin{definition}
	Im Fall $Y = \lr{1}$ gilt $x~R~x' \Leftrightarrow x = x' \vee x, x' \in A$. Dann heißt $Y \cup_f X$ \tib{Zusammenschlagen von $A$} und wird auch $X / A$ geschrieben.
\end{definition}
\begin{proposition}
	Unter geeigneten Zusatzvorraussetzungen gilt:
	\begin{enumerate}
		\item Die Inklusion $i : Y \to Y \cup_f X$ ist eine stetige injektive Einbettung.
		\item Die Inklusion $X \setminus A \to Y \cup_f X$ ist eine abgeschlossene Einbettung.
	\end{enumerate}
\end{proposition}
\begin{definition}
	Die folgenden Konstruktionen bilden das Möbiusband:
	\begin{enumerate}
		\item Es seien $X = [0,1)_\bR \times [0,1]_\bR$ und $Y = \lr{1} \times [0,1]_\bR$. Es sei
		\begin{align*}
			A = \lr{0} \times [0,1]_\bR.
		\end{align*}
		So verklebt die Funktion
		\begin{align*}
			f : A &\to Y\\
			(0,y) &\mapsto (1,1-y)
		\end{align*}
		die Ränder $A$ und $Y$ "verdreht".
		\item Es sei $X = [0,1]_\bR \times [0,1]_\bR$ und $Y = [0,1]_\bR$. Es sei
		\begin{align*}
			A = \lr{(x,z) \in X : x = 0 \vee x = 1}.
		\end{align*}
		So verklebt die Funktion
		\begin{align*}
			f : A &\to Y\\
			(0,y) &\mapsto 1-y\\
			(1,y) &\mapsto y
		\end{align*}
		analog die Ränder
	\end{enumerate}
\end{definition}
\chapter{Zusammenhang}
\begin{definition}
	Ein topologischer Raum $(X, \cO)$ heißt \tib{zusammenhängend}, wenn es keine nichtleeren, disjunkten, offenen Mengen $U_1, U_2$ gibt, sodass $U_1 \cap U_2 = X$.
\end{definition}
\begin{lemma}
	$(X, \cO)$ ist genau dann zusammenhängend, wenn $\emptyset$ und $X$ die einzigen geschloffenen Teilmengen von $X$ sind.
\end{lemma}
\begin{theorem}
	Jedes offene Intervall $(a,b) \subset \bR$ ist zusammenhängend.
\end{theorem}
\begin{proof}
	Angenommen, $(a,b)$ sei nicht zusammenhängend. Es seien $O_1, O_2$ nichtleere disjunkte offene Mengen mit $O_1 \cup O_2 = (a,b)$. OBdA gebe es $x_1 \in O_1$, $x_2 \in O_2$, sodass $x_1 < x_2$. Wir setzten
	\begin{align*}
		S := \lr{s \in (a,b) \mid [x_1, s] \subseteq O_1}
	\end{align*}
	So ist $S$ durch $x_2$ nach oben beschränkt. Da $\bR$ vollständig ist, hat $S$ ein reelles Supremum $m$. 
	\begin{enumerate}
		\item Angenommen, $m \in O_1$. Dann existiert ein $\epsilon > 0$, sodass $U_\epsilon(m) \subseteq O_1$. Dann ist aber $m + \frac{\epsilon}{2} \in O_1$, also ist $m$ keine obere Schranke von $S$.
		\item Angenommen, $m \in O_2$. Dann existiert ein $\epsilon \in > 0$, sodass $U_\epsilon(m) \subseteq O_2$. Dann ist aber $m - \frac{\epsilon}{2} \in O_2$, also $m - \frac{\epsilon}{2} \notin S$, also existiert eine kleinere obere Schranke von $S$.
	\end{enumerate}
	Da beide Fälle zum Widerspruch führen muss $(a,b)$ zusammenhängend sein.
\end{proof}
\begin{corollary}
	$\bQ = (-\infty, \sqrt{2}) \cup (\sqrt{2}, \infty)$ ist als Teilmenge von $\bR$ nicht zusammenhängend.
\end{corollary}
\pagebreak
\begin{theorem}
	Sei $(X, \cO)$ ein topologischer Raum, $A \subseteq X$ zusammenhängend und $B$ beliebig. Dann gilt:
	\begin{enumerate}
		\item Wenn $A \subseteq B \subseteq \ol A$ gilt, ist $B$ zusammenhängend.
		\item Wenn $A$ sowohl Punkte aus dem Inneren als auch aus dem Äußeren von $B$ enthält, so enthält $A$ auch Punkte aus $\partial B$.
	\end{enumerate}
\end{theorem}
\begin{theorem}
	Es seien $(X, \cO_X)$, $(Y, \cO_Y)$ topologische Räume. Es sei $X$ zusammenhängend und $f : X \to Y$ stetig. Dann ist $f[X]$ zusammenhängend.
\end{theorem}
\begin{lemma}
	Ein topologischer Raum $(X, \cO)$ ist genau dann zusammenhängend, wenn es zu jeder offenen Überdeckung $U$ von $X$ und für alle $a,b \in X$ eine endliche Teilüberdeckung $U_1, \hdots, U_n$ gibt, sodass:
	\begin{enumerate}
		\item $U_1$ ist das einzige $U_i$, welches $a$ enthält: 
		$$\forall j > 1: a \in U_1 \setminus U_j$$
		\item $U_n$ ist das einzige $U_i$, welches $b$ enthält: 
		$$\forall j < n: b \in U_n \setminus U_j$$
		\item Jedes $U_i$ überschneidet sich nur mit $U_{i-1}$, $U_i$ und $U_{i+1}$:
		$$U_i \cap U_j \Leftrightarrow \abs{i-j} \leq 1$$
	\end{enumerate}
\end{lemma}
\begin{proof}
	\begin{enumerate}
		\item[$\Leftarrow$] Sei $X$ nicht zusammenhängend. Dann gilt $X = U_1 \cup U_2$ mit $U_1, U_2$ disjunkt, also gibt es für $a \in U_1$, $b \in U_2$, $U = \lr{U_1, U_2}$ keine Möglichkeit, Eigenschaft $3$ zu erfüllen.
		\item
	\end{enumerate}
\end{proof}
\begin{theorem}
	Es sei $(X, \cO)$ ein topologischer Raum. Es seien $A,B \subseteq X$ zusammenhängendund $A \cap B \neq \emptyset$. Dann ist $A \cup B$ zusammenhängend.
\end{theorem}
\begin{theorem}
	Es seien $I$ eine Indexmenge und $(X_i, \cO_i)_{i \in I}$ topologische Räume. Dann ist der Produktraum $\prod_{i \in I} X_i$ genau dann zusammenhängend, wenn alle $(X_i, \cO_i)$ zusammenhängend sind.
\end{theorem}
\begin{proof}
	(...)\\
	$Z_j$ ist zusammenhängend, da $X_j$ zusammenhängend ist und jede disjunkte Zerlegung von $Z_j$ in Teilmengen auch $X_j$ zerlegen müsste.
	(...)
\end{proof}
\begin{theorem}
	Sei $(X, \cO)$ ein topologischer Raum.
	Die \tib{Zusammenhangskomponente} eines Punkts $x \in X$ ist
	\begin{align*}
		K(x) = \bigcup \lr{Y \subseteq X : x \in Y, Y \tn{ zusammenhängend}}
	\end{align*}
\end{theorem}
\begin{theorem}
	Ein topologischer Raum $(X, \cO)$ heißt \tib{total unzusammenhängend}, falls $\forall x \in X : K(x) = \lr{x}$.
\end{theorem}
\begin{definition}
	Sei $(X, \cO)$ ein topologischer Raum. Es seien $x, y \in X$. Ein \tib{Weg von $x$ nach $y$} ist eine stetige Abbildung $\gamma : [0,1]_\bR \to X$, sodass $\gamma(0) = x$ und $\gamma(1) = y$.
	Eine Abbildung $\gamma : [0,1]_\bR \to X$ 
\end{definition}
\begin{definition}
	Ein topologischer Raum $(X, \cO)$ heißt \tib{wegzusammenhängend}, falls es für jedes paar $x, y \in X$ einen Weg von $x$ nach $y$ gibt. $K_w(x)$ ist die \tib{Wegzusammenhangskomponente} von $X$.
\end{definition}
\begin{definition}
	Ein topologischer Raum $(X, \cO)$ heißt \tib{lokal zusammenhängend,} falls für jeden Punkt $x \in X$ und jede Umgebung $U \in U(x)$ eine Umgebung $V \in U(x)$ gibt, sodass $V \subseteq U$ und sodass $V$ zusammenhängend ist. Ein topologischer Raum $(X, \cO)$ heißt \tib{lokal wegzusammenhängend}, falls ein solches $V$ existiert, welches wegzusammenhngend ist.
\end{definition}
\begin{theorem}
	Jeder wegzusammenhängende Raum ist zusammenhängend.
\end{theorem}
\begin{proof}
	Es sei $X$ nicht zusammenhängend und $X = U \cup V$ mit $U,V$ disunkt, nichtleer und offen. Angenommen, es gäbe einen Weg $\gamma$ von $x \in U$ nach $y \in V$. So gilt $[0,1]_\bR = \gamma^{-1}(U) \cup \gamma^{-1}(V)$, aber da $\gamma$ stetig ist sind $\gamma^{-1}(U)$, $\gamma^{-1}(V)$ nichtleer und offen, und da $U$ und $V$ disjunkt sind, sind auch die Urbilder disjunkt. Somit wäre $[0,1]_\bR$ nicht zusammenhängend. Widerspruch!!!
\end{proof}
\begin{theorem}
	Wenn $X$ lokal wegzusammenhängend ist, ist $K_w(x)$ offen für alle $x$.
\end{theorem}
\begin{theorem}
	Wenn $X$ lokal zusammenhängend ist, ist $K(x)$ offen für alle $x$.
\end{theorem}
\begin{theorem}
	$\prod_{j \in J} X_j$ ist lokal zusammenhängend genau dann, wenn jedes $X_i$ lokal zusammenhängend ist und alle bis auf eindlich viele $X_i$ zusammenhängend sind.
\end{theorem}
\begin{theorem}
	Der Graph des topologischen Sinus ist zusammenhängend, aber nicht lokal wegzusammenhängend.
\end{theorem}
\begin{theorem}
	Der topologische Besen ist wegzusammenhängend, aber nicht lokal zusammenhängend.
\end{theorem}
\chapter{Filter und Konvergenz}
\begin{definition}
	Es sei $(X, \cO)$ ein topologischer Raum. Eine Folge $(x_n)_{n \in \bN}$ in $X$ \tib{konvergiert gegen $x$}, wir schreiben $x_n \to x$, wenn
	\begin{align*}
		\forall U \in U(x) : \exists n_0 \in \bN : \forall n \geq n_0 ; x_n \in U
	\end{align*}
\end{definition}
\begin{definition}
	Ein Punkt $x \in X$ heißt \tib{Häufungspunkt einer Folge} $(x_n)_{n \in \bN}$, wenn
	\begin{align*}
		\forall U \in U(x) : \forall n_0 \in \bN : \exists n \geq n_0 : x_n \in U
	\end{align*}
\end{definition}
\begin{definition}
	Ein Punkt $x \in X$ heißt \tib{Berührungspunkt einer Menge} $A \subseteq X$, wenn
	\begin{align*}
		\forall U \in U(x) : U \cap A \neq \emptyset
	\end{align*}
\end{definition}
\begin{theorem}
	Es sei $(X, \cO_X)$ ein topologischer Raum mit abzählbarer Umgebungsbasis, und sei $x \in X$. Dann gilt $x \in \ol A$ genau dann, wenn es eine Folge in $A$ gibt, die gegen $x$ konvergiert.
\end{theorem}
\begin{theorem}
	Es seien $(X, \cO_X)$ und $(Y, \cO_Y)$ topologische Räume mit abzählbarer Umgebungsbasis, und sei $x \in X$. Dann ist $f : X \to Y$ stetig in $x$ genau dann, wenn
	\begin{align*}
		\forall (x_n)_{n \in \bN} : x_n \to x \implies f(x_n) \to f(x)
	\end{align*}
\end{theorem}
\begin{definition}
	Eine Teilmenge $\cF \subseteq \cP(X)$ heißt \tib{Filter über $X$}, wenn:
	\begin{enumerate}
		\item $X \in \cF$, $\emptyset \notin \cF$.
		\item Ist $A \in \cF$ und $A \subseteq B$, so ist $B \in \cF$.
		\item Sind $A$ und $B$ in $\cF$, so auch $A \cap B$.
	\end{enumerate}
\end{definition}	
\begin{theorem}
	Es sei $(X, \cO)$ ein topologischer Raum und $x \in X$. Dann ist die Menge $U(x)$ der Umgebungen von $x$ ein Filter über $X$, genannt der \tib{Umgebungsfilter}.
\end{theorem}
\begin{definition}
	Ein Filter $\cF$ heißt \tib{frei} (en: non-principal), wenn $\bigcap F = \emptyset$. Falls $\cF$ nicht frei ist, so heißt $\cF$ \tib{fixiert}.
\end{definition}
\begin{theorem}
	Der \tib{Fréchet-Filter} 
	\begin{align*}
		\cF = \lr{A \subseteq \bN : \bN \setminus A \tn{ ist endlich}}
	\end{align*}
	ist ein freier Filter über $\bN$.
\end{theorem}
\begin{definition}
	Es sei $\cF$ ein Filter über $X$. So heißt eine Teilmenge $\cB \subseteq \cF$ eine \tib{Filterbasis von $\cF$}, wenn
	\begin{align*}
		\forall F \in \cF : \exists B \in \cB : B \subseteq F
	\end{align*}
\end{definition}
\begin{example}
	Jede Umgebungsbasis eines Punkts ist gleichzeitig eine Filterbasis des Umgebungsfilters, und umgekehrt.
\end{example}
\begin{theorem}
	$\cB \subseteq \cP(X)$ ist Basis eines Filters aus $X$ genau dann, wenn $\emptyset \neq \cB$ und
	\begin{align*}
		\forall B_1, B_2 \in \cB: \exists B_3 \in \cB : B_3 \subseteq B_1 \cap B_2
	\end{align*}
\end{theorem}
\begin{definition}
	Ein Filter $U \subseteq \cP(X)$ heißt \tib{Ultrafilter}, falls $U$ bezüglich der Inklusion maximal ist.
\end{definition}
\begin{theorem}
	\theoremname{Lemma von Zorn:} Jede induktive Halbordnung hat ein maximales Element.
\end{theorem}
\begin{corollary}
	Jeder Filter kann zu einem Ultrafilter erweitert werden.
\end{corollary}
\begin{definition}
	Eine Halbordnung heißt \tib{induktiv}, falls jede Kette eine obere Schranke hat.
\end{definition}
(...)
\chapter{Trennungseigenschaften}
\begin{definition}
	Ein topologischer Raum $(X, \cO)$ heißt:
	\begin{enumerate}
		\item \tib{T0}, falls je zwei Punkte unterschiedliche Umgebungsfilter haben:
		\begin{align*}
			\forall x \neq y : U(x) \neq U(y)
		\end{align*}
		\item \tib{T1}, falls je zwei Punkte Umgebungen haben, die jeweils nicht den anderen Punkt enthalten:
		\begin{align*}
			\forall x \neq y : \exists U_x \in U(x), U_y \in U(y) : x \notin U_y, y \notin U_x
		\end{align*}
		\item \tib{T2} oder \tib{Hausdorffsch}, falls je zwei Punkte disjunkte Umgebungen haben:
		\begin{align*}
			\forall x \neq y : \exists U_x \in U(x), U_y \in U(y) : U_x \cap U_y \neq \emptyset
		\end{align*}
		\item \tib{T3}, falls jede abgeschlossene Menge zu jedem Punkt außerhalb der Menge disjunkte Umgebungen haben:
		\item \tib{T3a} (oder \tib{T3$\frac{1}{2}$}), falls jede abgeschlossene Menge von jedem Punkt außerhalb der Menge durch eine stetige Funktion getrennt werden kann.
		\item \tib{T4}, falls je zwei disjunkte geschlossene Mengen disjunkte Umgebungen haben.
		\item \tib{Regulär}, falls er sowohl $T1$ als auch $T3$ ist,
		\item \tib{Vollständig regulär,} falls er sowohl $T1$ als auch $T3a$ ist,
		\item \tib{Normal}, falls er sowohl $T1$ als auch $T4$ ist.
	\end{enumerate}
\end{definition}
\begin{theorem}
	Die folgenden Aussagen sind äquivalent:
	\begin{enumerate}
		\item $(X, \cO)$ ist $T1$,
		\item Alle einelementigen Teilmengen von $X$ sind abgeschlossen,
		\item Jedes $A \subseteq X$ ist der Schnitt seiner offenen Obermengen.
	\end{enumerate}
\end{theorem}
\begin{theorem}
	Die folgenden Aussagen sind äquivalent:
	\begin{enumerate}
		\item $(X, \cO)$ ist Hausdorff,
		\item Die Diagonalmenge
		\begin{align*}
			\Delta = \lr{(x,x) : x \in X}
		\end{align*}
		ist in der Produkttopologie abgeschlossen,
		\item Jedes $x \in X$ ist der Schnitt der Abschlüsse seiner Umgebungen,
		\item Jeder konvergente Filter hat genau einen Berührpunkt.
	\end{enumerate}
\end{theorem}
\begin{theorem}
	Sei $(X, \cO)$ ein vollständig regulärer Raum. Dann ist die Abbildung $e$, definiert durch
	\begin{align*}
		L &= \lr{f : f : X \to [0,1] \tn{stetig}}\\
		e : X &\to \prod_{f \in L} [0,1]\\
		x &\mapsto (f(x))_{f \in L},
	\end{align*}
	eine Einbettung in einen Produktraum $\prod_{f \in L} [0,1]$.
\end{theorem}
\begin{theorem}
	Jeder Unterraum eines $T0$-Raumes ist $T0$.
\end{theorem}
\begin{theorem}
	Jeder Unterraum eines $T1$-Raumes ist $T1$.
\end{theorem}
\begin{theorem}
	Jeder Unterraum eines $T2$-Raumes ist $T2$.
\end{theorem}
\begin{theorem}
	Jeder Unterraum eines $T3$-Raumes ist $T3$.
\end{theorem}
\begin{theorem}
	Jeder Unterraum eines $T3a$-Raumes ist $T3a$.
\end{theorem}
\begin{corollary}
	Es seien $X_i$, $i \in I$ nichtleere topologische Räume. Dann ist der Produktraum $\prod_{i \in I} X_i$ genau dann $T0/T1/T2/T3/T3a$, wenn jedes $X_i$ $T0/T1/T2/T3/T3a$ ist.
\end{corollary}
\begin{theorem}
	Jeder \ul{abgeschlossene} Unterraum eines $T4$-Raumes ist $T4$.
\end{theorem}
\begin{theorem}
	Die $T4$-Eigenschaft ist nicht gegen Produkte abgeschlossen:
	Die Sorgenfreygerade $\bR$ mit der durch $[a,b)$ erzeugten Topologie ist normal, ihr Produkt mit sich selbst aber nicht.
\end{theorem}
\end{document}